\subsection{\texorpdfstring{例 I：\(A_{n}\) 型紧李代数}{例 I：A\_\{n\} 型紧李代数}}

设 V 是有限维复向量空间，\(\Phi\) 是 V 上分离的正埃尔米特型。与 \(\Phi\) 相伴的酉群（参见《代数学》，第 IX 章）是子群 \(\mathbf{U}(\Phi)\) （\(\mathbf{GL}(V)\) 中由复希尔伯特空间 \((V,\Phi)\) 的自同构组成的）；这是群 \(\mathbf{GL}(V)\) 的一个（实）李子群，其李代数是子代数 \(\mathfrak{u}(\Phi)\) （实李代数 \(\mathfrak{gl}(V)\) 中由满足 \(x^{*} = -x\) 的 V 的自同态 x 组成的）（第 III 章，§3，第 10 节，命题 37 的推论 2），其中 \(x^{*}\) 表示 x 相对于 \(\Phi\) 的伴随。由于群 \(\mathbf{U}(\Phi)\) 紧（§1 第 1 节），\(\mathfrak{u}(\Phi)\) 是紧实李代数。同样，特殊酉群 \(\mathbf{SU}(\Phi) = \mathbf{U}(\Phi) \cap \mathbf{SL}(V)\) 是 \(\mathbf{SL}(V)\) 的紧李子群，其李代数为 \(\mathfrak{su}(\Phi) = \mathfrak{u}(\Phi) \cap \mathfrak{sl}(V)\) 。

当 \(V = C^{n}\) 且 \(\varPhi\) 是通常的埃尔米特型（\(C^{n}\) 的典则基关于它标准正交）时，我们把 \(\mathbf{U}(n, \mathbf{C})\) 、\(\mathbf{SU}(n, \mathbf{C})\) 、\(\mathfrak{u}(n, \mathbf{C})\) 、\(\mathfrak{s}\mathfrak{u}(n, \mathbf{C})\) 分别写作 \(\mathbf{U}(\varPhi)\) 、\(\mathbf{SU}(\varPhi)\) 、\(\mathfrak{u}(\varPhi)\) 、\(\mathfrak{s}\mathfrak{u}(\varPhi)\) 。\(\mathbf{U}(n, \mathbf{C})\)（相应地，\(\mathfrak{u}(n, \mathbf{C})\)）的元素

是矩阵 \(A \in \mathrm{M}_{n}(\mathbf{C})\) （满足 \(A.^{t}\bar{A} = I_{n}\) ，相应地 \(A = -^{t}\bar{A}\) ），分别称为酉的（相应地，反埃尔米特的）。

命题 4. a) 复李代数 \(\mathfrak{sl}(V)\) 的紧实形式是诸代数 \(\mathfrak{su}(\Phi)\) ，其中 \(\Phi\) 取遍复向量空间 \(V\) 上分离的正埃尔米特型所成的集。

\begin{enumerate}
\def\labelenumi{\alph{enumi})}
\setcounter{enumi}{1}
\tightlist
\item
  诸代数 \(\mathfrak{u}(\varPhi)\) 是 \(\mathfrak{gl}(\mathrm{V})\) 的紧实形式。
\end{enumerate}

设 \(\varPhi\) 是 V 上分离的正埃尔米特型。对一切 \(x\in\mathfrak{gl}(V)\) ，令 \(\sigma(x)=-x^{*}\) （其中 \(x^{*}\) 是 \(x\) 相对于 \(\varPhi\) 的伴随）。则 \(\sigma\) 满足第 1 节命题 1 的条件 (2)，故不动点集 \(\mathfrak{u}(\varPhi)\) （相应地，\(\mathfrak{su}(\varPhi)\) ）——即 \(\sigma\) 在 \(\mathfrak{gl}(V)\) （相应地，\(\mathfrak{sl}(V)\) ）上的不动点所成的集——是 \(\mathfrak{gl}(V)\) （相应地，\(\mathfrak{sl}(V)\) ）的紧实形式。由于 \(\mathbf{GL}(V)\) 传递地作用在 V 上分离的正埃尔米特型所成的集上（《代数学》，第 IX 章），也传递地作用在 \(\mathfrak{sl}(V)\) 的紧实形式所成的集上（第 3 节，定理 1，及第 VIII 章，§13，第 1 节 (VII)），命题 4 得证。

推论. \(\mathrm{A}_n\) 型（ \(n \geq 1\) ）的每个紧单实李代数同构于 \(\mathfrak{su}(n + 1, \mathbf{C})\) 。

事实上，每个 \(\mathbf{A}_n\) 型复李代数都同构于 \(\mathfrak{sl}(n + 1,\mathbf{C})\) （第 VIII 章，§13，第 1 节）。

注。1) 我们有 \(\mathfrak{gl}(\mathrm{V}) = \mathfrak{sl}(\mathrm{V}) \times \mathbf{C}.1_{\mathrm{V}}\) ，\(\mathfrak{u}(\varPhi) = \mathfrak{su}(\varPhi) \times \mathbf{R}.i1_{\mathrm{V}}\) ；\(\mathfrak{gl}(\mathrm{V})\) 的紧实形式是诸 \(\mathfrak{su}(\varPhi) \times \mathbf{R}. \alpha 1_{\mathrm{V}}, \alpha \in \mathbf{C}^*\) 。

\begin{enumerate}
\def\labelenumi{\arabic{enumi})}
\setcounter{enumi}{1}
\tightlist
\item
  若复李代数 \(\mathfrak{a} = \mathfrak{sl}(n,\mathbf{C})\) 装备了第 VIII 章 §13 第 1 节 (IX) 中引入的分裂与谢瓦莱系，则采用第 2 节的记号，
\end{enumerate}

\[
\mathfrak {a} _ {u} = \mathfrak {s u} (n, \mathbf {C}), \quad \mathfrak {a} _ {0} = \mathfrak {s l} (n, \mathbf {R}), \quad \mathfrak {a} _ {u} \cap \mathfrak {a} _ {0} = \mathfrak {o} (n, \mathbf {R}).
\]

